Throughout, as in the proof of Theorem \claimref{synthesis.strict-completions-exist-exactly-when-the-signature-does-not-jump}: where a continuous family of symmetric maps is invertible, its spectral projections onto the positive and negative parts are continuous, so $G_{\pm}$ are vector bundles over $\partial B$ and $\operatorname{sgn} g$, the difference of the two projections, is continuous there.

\emph{(1).} By clause (4) of the cited Theorem, the base $\R^{m}$ being contractible, $G_{\pm} \oplus (K_{0})_{\pm}$ are trivial on $\R^{m} \setminus B$, which contains $\partial B$; $(K_{0})_{+}$ and $(K_{0})_{-}$ are fixed subspaces of dimensions $c$ and $c'$.

\emph{(2).} Put $A = \R^{n} \oplus \R^{c} \oplus \R^{c'}$, of dimension $N$, and for $y \in \partial B$ let $P(y) = G_{+}(y) \oplus \R^{c} \oplus 0$ and $Q(y) = G_{-}(y) \oplus 0 \oplus \R^{c'}$: orthogonal complements of each other in $A$, of dimensions $a = p + c$ and $b = q + c'$. By hypothesis each has a continuous frame; the Gram-Schmidt process, which is continuous, makes the frames orthonormal, and $F(y)$, with these $a + b = N$ vectors as columns, is a continuous map $\partial B \to \mathrm{O}(N)$ with $F\,\operatorname{diag}(1_{a}, -1_{b})\,F^{-1} = \pi_{P} - \pi_{Q} = \operatorname{sgn} g \oplus 1_{c} \oplus (-1_{c'})$. The fibre of the completion is $A \oplus \R^{N}$, the second summand being the third summand of $W_{0}$, on which $K_{0}$ is the identity; in this order $K_{0}$ and $\operatorname{sgn} g$ together are $F\operatorname{diag}(1_{a}, -1_{b})F^{-1} \oplus 1_{N}$.

\emph{The rotation.} In blocks for $A \oplus \R^{N}$, $R_{t} = \begin{psmallmatrix} \cos t & -\sin t \\ \sin t & \cos t\end{psmallmatrix}$ with scalar blocks. It is orthogonal, $R_{0} = 1$, and $R_{\pi/2}\begin{psmallmatrix} 1 & 0 \\ 0 & X\end{psmallmatrix}R_{\pi/2}^{-1} = \begin{psmallmatrix} X & 0 \\ 0 & 1\end{psmallmatrix}$ for every $X$. So $\Phi_{t}$ is a product of orthogonal maps, continuous in $(t, y)$, $\Phi_{0} = F \oplus F^{-1}$ and $\Phi_{\pi/2} = (F \oplus 1)(F^{-1} \oplus 1) = 1$.

\emph{The form.} Each line of the definition of $\hat h$ is continuous in $(r, \omega)$ and symmetric ($\Phi_{t}^{-1} = \Phi_{t}^{\mathsf T}$). The lines agree where they meet. At $r = \rho$: $\lambda = 0$ and $\hat y = y$. At $r = \tfrac23\rho$: $\lambda = 1$ gives $\operatorname{sgn} g(\hat y) \oplus K_{0}$, and $t = 0$ gives $\Phi_{0}D\Phi_{0}^{-1} = F\operatorname{diag}(1_{a}, -1_{b})F^{-1} \oplus F^{-1}1_{N}F$, the same map by the previous paragraph. At $r = \tfrac13\rho$: $t = \tfrac\pi2$ gives $D$. Near $r = 0$ the map is the constant $D$, so it is continuous at the centre, where polar coordinates are singular. It is invertible at every point: on $\R^{m} \setminus B$ both summands are; on the second line $(1-\lambda)g + \lambda\operatorname{sgn} g$ has the eigenvectors of $g(\hat y)$ with eigenvalues $(1-\lambda)\mu + \lambda\operatorname{sgn}\mu$, of the sign of $\mu$; on the third and fourth it is an orthogonal conjugate of $D$. It equals $g \oplus K_{0}$ on $\R^{m} \setminus B$. For $m = 1$ the sphere is two points and nothing changes. \emph{Other forms of the same signature.} If $K_{0}'$ is any nondegenerate form on $W_{0}$ of signature $(c + N, c')$, then by Sylvester's law of inertia $K_{0}' = \Lambda^{\mathsf T}K_{0}\Lambda$ for a constant invertible $\Lambda$ of $W_{0}$, and $(1 \oplus \Lambda)^{\mathsf T}\,\hat h\,(1 \oplus \Lambda)$ is continuous, symmetric, invertible at every point and equal to $g \oplus K_{0}'$ on $\R^{m} \setminus B$. At a point of $\Sig$ the restriction of $\hat h$ to $\R^{n}$ is not $g$ in general (at the centre it is the upper left block of $D$).

\emph{(3).} $G_{+} \oplus G_{-} = \R^{n}$ over $\partial B$. If $G_{+} \oplus \R^{c}$ is trivial, of rank $p + c$, then $G_{-} \oplus \R^{p + c} \cong G_{-} \oplus G_{+} \oplus \R^{c} = \R^{n + c}$ is trivial, and symmetrically. With (1) and (2) this gives the equivalence.

\emph{(4), the lifting clause.} We use clause (L) of Theorem \claimref{synthesis.position-vector-eigenbundle-in-codimension-three-has-a-class-of-order-two}: for $n \geq 1$ and a continuous $\kappa \colon [0,1] \times [0,1] \to \mathrm{SO}(n)$ with $\kappa(u,0) = \kappa(u,1)$ for every $u$ and $\kappa(1,t)$ independent of $t$, there is a continuous map $\hat\kappa$ of the square into the rotors of $\Cl(n,0)$ over $\kappa$, with $\hat\kappa(u,0) = \hat\kappa(u,1)$ and $\hat\kappa(1,t)$ independent of $t$, and every continuous path of rotors over $\kappa(0,\cdot)$ is $\pm\hat\kappa(0,\cdot)$, so it is closed. It is proved there from Theorem \claimref{std-algebra.spin-double-cover}; the rotors are the even versors of unit vectors.

\emph{(4), no completion of rank two.} The eigenvectors: $G(x,y)^{\mathsf T}$ applied to $(x,y)$ gives $(x^{2} + y^{2})(x,y)$, and to $(-y,x)$ gives $-(x^{2} + y^{2})(-y,x)$; both lines have these nonvanishing sections on every circle about the origin. Suppose $h$ is a continuous symmetric form on $\R^{2}$ over the plane, nondegenerate at every point and equal to $G$ outside a bounded set; let $r$ be a radius beyond that set. Its signature is constant, so $(1,1)$, and $\operatorname{sgn} h$ is continuous, a reflection at every point. Let $s_{0} = \operatorname{diag}(1,-1)$ and $\kappa_{0}(y) = \operatorname{sgn} h(y)\,s_{0}$, a product of two reflections, so a continuous map of the plane into $\mathrm{SO}(2)$. At the point of the circle of radius $r$ with angle $a$, $\operatorname{sgn} h = \operatorname{sgn} G = \begin{psmallmatrix} \cos 2a & \sin 2a \\ \sin 2a & -\cos 2a\end{psmallmatrix}$ and $\kappa_{0}$ is the rotation through $2a$. Put $\kappa(u,t) = \kappa_{0}\bigl((1-u)\,r\,(\cos 2\pi t, \sin 2\pi t)\bigr)$: it satisfies the hypotheses of (L) with $n = 2$, and $\kappa(0,t)$ is the rotation through $4\pi t$. With $f_{1}, f_{2}$ the basis of the plane, the rotors of $\Cl(2,0)$ are the elements $\cos b + \sin b\,f_{1}f_{2} = f_{1}(\cos b\,f_{1} + \sin b\,f_{2})$ (a product of two unit vectors of the plane is of this form, and a product of such elements is again one), and $\cos b + \sin b\,f_{1}f_{2}$ induces the rotation through $-2b$ ($f_{1} \mapsto \cos 2b\,f_{1} - \sin 2b\,f_{2}$). So the rotor with $b = -2\pi t$ is a lift of $\kappa(0,\cdot)$, and by (L) $\hat\kappa(0,t)$ is that rotor or its negative, the rotor with $b = -2\pi t + \pi$. The map $j \colon \cos b + \sin b\,f_{1}f_{2} \mapsto$ (rotation through $b$) is a continuous isomorphism of the group of rotors of the plane onto $\mathrm{SO}(2)$. So $j \circ \hat\kappa$ again satisfies the hypotheses of (L), and $j(\hat\kappa(0,t))$ is the rotation through $-2\pi t + b_{0}$ with $b_{0} = 0$ or $\pi$. A lift of it is the rotor with $b = \pi t - \tfrac12 b_{0}$, which at $t = 1$ is the negative of its value at $t = 0$. This contradicts (L). (For $\begin{psmallmatrix} x & y \\ y & -x\end{psmallmatrix}$ the same argument stops at the first step: the loop is the rotation through $2\pi t$, whose lift does not close.)

\emph{(4), one added direction.} $Q(\nu) = \beta^{2}(x^{2} + y^{2}) + (1 - \beta)^{2}$ is positive: where $\beta = 1$ the point is not the origin. So $\nu$ is a continuous nowhere-vanishing vector field in $\R^{3}$ and $1 - 2\nu\nu^{\mathsf T}/Q(\nu)$ is the reflection along it, continuous, symmetric, of signature $(2,1)$. Where $\beta = 0$ it is $\operatorname{diag}(1,1,-1)$. Where $\beta = 1$, $\nu = (-y, x, 0)$ and the reflection is $+1$ on $(x,y,0)$ and on the third direction and $-1$ on $(-y,x,0)$: it is $\operatorname{sgn} G \oplus 1$. Further out, $\bigl((1-\lambda) + \lambda(x^{2} + y^{2})\bigr)\operatorname{sgn} G \oplus 1$ with $\lambda$ rising from $0$ to $1$ joins it to $G \oplus 1 = (x^{2} + y^{2})\operatorname{sgn} G \oplus 1$ through nondegenerate forms.

\emph{(5).} For one copy, clause (5) of the cited Theorem. For two copies on $\R^{2} \oplus \R^{2}$, with $\ell = (C,S)$, $\ell^{\perp} = (-S,C)$ the unit eigenvectors of one copy at the point of angle $a$, $C = \cos\tfrac a2$, $S = \sin\tfrac a2$: the vectors $(C\ell, S\ell)$, $(-S\ell, C\ell)$ are orthonormal and span $\ell \otimes \R^{2}$, the positive eigenbundle, and $(C\ell^{\perp}, S\ell^{\perp})$, $(-S\ell^{\perp}, C\ell^{\perp})$ the negative one; their components are quadratic in $(C,S)$, so unchanged under $(C,S) \mapsto (-C,-S)$ and continuous functions of the point. (2) applies with $c = c' = 0$, $N = 4$. For the families of clause (6) of the cited Theorem: over $\R$ the two points of the sphere carry eigenspaces of equal dimension, which is triviality over two points; over $\R^{2}$ the form $x^{2} + y^{2}$ on a line has the constant positive line.
